\documentclass[manuscript,review,anonymous]{acmart}
% PerceptFence anonymous review manuscript skeleton.
% IUI 2027 CFP re-checked 2026-04-26:
% https://iui.acm.org/2027/call-for-papers/
% Keep author names, affiliations, acknowledgements, organization names,
% public repository links, and non-synthetic data out of this file until
% the camera-ready swap described in paper/authors.private.md.

\setcopyright{none}
\settopmatter{printacmref=false}
\renewcommand\footnotetextcopyrightpermission[1]{}

\acmConference[IUI '27]{The Annual ACM Conference on Intelligent User Interfaces}{February 8--11, 2027}{Helsinki, Finland}
\acmYear{2027}
\copyrightyear{2027}
\acmDOI{}
\acmISBN{}
\acmSubmissionID{anonymous-review-draft}

\title[Consent-Aware Runtime Mediation for Screen-Share AI]{Consent-Aware Runtime Mediation for Privacy in Real-Time Screen-Share AI Assistants}

\author{Anonymous Author(s)}
\affiliation{%
  \institution{Anonymous Institution(s)}
  \country{Anonymous Country}}
\email{anonymous@example.invalid}

\begin{abstract}
Live screen-share AI assistants can observe raw screen and speech streams, but users lack fine-grained runtime control over what the assistant may see, retain, and say.
Existing chatbot and privacy controls do not address this control gap because the assistant must use fast-changing interface context without exposing, remembering, or repeating information outside current user consent.
We study a runtime mediation layer that enforces consent, redaction, memory controls, and output safeguards for live multimodal assistance between capture adapters, session memory, and assistant responses.
We will evaluate the layer on synthetic screen-share scenarios covering secrets, personal data, sensitive speech fragments, prompt-injection text, dense interfaces, zoom changes, and fast window switching, measuring sensitive exposure, false blocks, latency overhead, and task completion against an unmediated prototype baseline.
The planned contribution is a bounded design and evaluation of runtime control mechanisms for screen-share assistants, with claims limited to synthetic evidence until benchmark results are available.
\end{abstract}

\begin{CCSXML}
<ccs2012>
   <concept>
       <concept_id>10003120.10003121.10003124.10010866</concept_id>
       <concept_desc>Human-centered computing~User interface design</concept_desc>
       <concept_significance>500</concept_significance>
   </concept>
   <concept>
       <concept_id>10003120.10003121.10011748</concept_id>
       <concept_desc>Human-centered computing~Empirical studies in HCI</concept_desc>
       <concept_significance>300</concept_significance>
   </concept>
   <concept>
       <concept_id>10002978.10003022.10003465</concept_id>
       <concept_desc>Security and privacy~Human and societal aspects of security and privacy</concept_desc>
       <concept_significance>300</concept_significance>
   </concept>
</ccs2012>
\end{CCSXML}

\ccsdesc[500]{Human-centered computing~User interface design}
\ccsdesc[300]{Human-centered computing~Empirical studies in HCI}
\ccsdesc[300]{Security and privacy~Human and societal aspects of security and privacy}

\keywords{screen-share assistants, multimodal AI, runtime mediation, consent controls, redaction, session memory}

\begin{document}
\maketitle

\section{Introduction}
% TODO(the section-scope task): Insert the blinded section scope paragraph. State the
% user-control gap using the frozen thesis and motivate why real-time screen and
% speech streams require runtime mediation instead of static chatbot settings.
% Do not claim achieved reductions until eval/results/ contains benchmark outputs.

\section{Problem and Threat Model}
% Placeholder input for the Problem and Threat Model section.
% the threat-model task owns the fuller threat-model text. Keep this file anonymous and
% replace only with blinded, synthetic-scope content.

\subsection{Threat Model Scope}
% TODO(the threat-model task): Define the session model, data surfaces, user-control goals,
% and threat model. Scope to accidental exposure and model/tool behavior in
% synthetic sessions. Exclude OS compromise, malicious capture drivers, and
% real-user data.

\subsection{Control Surface}
% TODO(the threat-model task): Describe observe, retain, and say controls as runtime
% permissions over screen regions/windows, speech fragments, session memory,
% and assistant output.

\subsection{Non-goals}
% TODO(the threat-model task): Explicitly state that this draft does not claim formal
% guarantees, differential privacy, side-channel resistance, or production
% readiness.


\section{System Design}
% TODO(the section-scope task): Insert the blinded section scope paragraph. Describe the
% mediation layer between capture inputs and assistant context/output, with
% explicit data-flow boundaries and no production-readiness claims.
\begin{figure}[t]
  \centering
  \fbox{\parbox{0.86\linewidth}{\centering Placeholder architecture figure: capture adapters $\rightarrow$ consent/policy engine $\rightarrow$ redaction engine $\rightarrow$ session memory gate $\rightarrow$ output guard $\rightarrow$ audit logger.}}
  \caption{Planned PerceptFence runtime mediation architecture. Replace this placeholder with the generated blinded figure after the post-filing benchmark/figure gate opens.}
  \Description{Placeholder box listing the planned data flow through capture, policy, redaction, memory, output guard, and audit components.}
  \label{fig:architecture-placeholder}
\end{figure}

\subsection{Capture Adapters}
% TODO: Summarize the synthetic screen and speech event schema.

\subsection{Redaction Engine}
% TODO: Describe sensitive-token and region-level transformations backed by
% synthetic fixtures.

\subsection{Session Memory Gate}
% TODO: Describe what may enter, remain in, or be removed from session memory.

\subsection{Output Guard}
% TODO: Describe response filtering for restricted content and prompt-injection
% text visible on screen.

\subsection{Audit Logger}
% TODO: Describe policy-decision logging. Do not claim tamper evidence unless
% an implementation and test are added.

\section{Policy and Consent Model}
% TODO(the section-scope task): Insert the blinded section scope paragraph. Specify policy
% objects for what the assistant may see, retain, and say at runtime, including
% session defaults and per-window or per-region overrides. Do not claim legal
% informed-consent coverage.

\section{Implementation}
% TODO(the section-scope task): Insert the blinded section scope paragraph. Map claims to
% runnable files under src/. Current expected modules: capture adapter,
% redaction engine, consent/policy engine, session memory gate, output guard,
% audit logger, synthetic fixture loader, and smoke-test harness.

\section{Evaluation}
% TODO(the section-scope task): Insert the blinded section scope paragraph. Convert this
% section from plan to results only after eval/results/baseline_vs_guarded.csv
% exists and has been verified.

\subsection{Synthetic Scenarios}
% Planned synthetic-only scenario classes: terminal secrets, chat-style messages,
% browser personal data, spoken sensitive fragments, prompt-injection text,
% dense interfaces, zoom changes, small fonts, and fast window switching.

\subsection{Metrics}
% Planned metrics: sensitive exposure rate, false block rate, latency overhead,
% and task completion against an unmediated prototype baseline.

\begin{table}[t]
  \caption{Benchmark table placeholder. Populate only from eval/results/baseline_vs_guarded.csv after the post-filing gate opens.}
  \Description{Placeholder table with rows for baseline and guarded modes and columns for planned evaluation metrics.}
  \label{tab:benchmark-placeholder}
  \centering
  \begin{tabular}{lllll}
    \toprule
    Condition & Sensitive exposure & False blocks & Latency overhead & Task completion \\
    \midrule
    Baseline & planned & planned & planned & planned \\
    Guarded & planned & planned & planned & planned \\
    \bottomrule
  \end{tabular}
\end{table}

\section{Limitations and Ethics}
% TODO(the section-scope task): Insert the blinded section scope paragraph. Include the short
% reviewer-facing ethics note required by the CFP. Current scope is
% synthetic-only and contains no non-synthetic screen captures, live-session
% captures, production telemetry, or human-subject data.

\section{Related Work}
% TODO(the section-scope task): Insert the blinded section scope paragraph. Add anonymized,
% third-person citations covering intelligent user interfaces, multimodal
% assistants, runtime permissions, redaction, assistant memory controls, and
% prompt-injection defenses. Keep self-citations in third person during
% double-blind review.

\section{GenAI Usage Disclosure}
% IUI 2027 asks for this section right before references when GenAI is used
% beyond editing the authors' own text. Replace this placeholder with the true
% disclosure before submission.
This section will disclose any use of generative AI tools across research design, code and data generation, analysis, and writing. It is currently a placeholder and must be updated before submission.

% TODO: Add references.bib and uncomment once citations are inserted.
% \bibliographystyle{ACM-Reference-Format}
% \bibliography{references}

\end{document}
